\subsection{Photon absorption in an alkali atom}

The three lowest electronic states of an alkali atom can be labeled as follows :
\begin{align}
    &\text{Ground electronic state:}         && 5s\ ^2S_{1/2} \notag \\
    &\text{First excited electronic state:}  && 5p\ ^2P_{1/2} \notag \\
    &\text{Second excited electronic state:} && 5p\ ^2P_{3/2} \notag
\end{align}
An electric dipole transition can take place between the $S$ and $P$ states with the
selection rules :
\begin{equation}
    \Delta L = 0,\,\pm 1 \quad (\text{but not } L=0 \to L=0),
    \qquad \Delta S = 0,
    \qquad \Delta J = 0,\,\pm 1
    \label{eq:dipole_selection}
\end{equation}
Thus, transitions from the ground state to both excited states are allowed.
 
\subsubsection*{Beer-Lambert Absorption Law}
 
In the optical pumping experiment, we are primarily interested in the absorption of light
by a volume of gas. Assuming the incident light is resonant with one of the above
transitions, a fraction of it will be absorbed by the atoms. Once excited, atoms decay back
to the ground state by spontaneous emission; however, since this emission occurs equally in
all directions, only a negligible fraction is re-radiated into the outgoing beam.
 
The attenuation of the photon flux through the gas follows :
\begin{equation}
    I = I_0\,e^{-\sigma_0 \rho \ell}
    \label{eq:beer_lambert}
\end{equation}
where $I_0$ and $I$ are the incident and outgoing photon flux, $\rho$ is the gas density,
$\ell$ is the path length through the gas, and $\sigma_0$ is the maximum absorption
cross-section at the center of the atomic resonance. The absorption coefficient $k_0$ is
related to $\sigma_0$ by :
\begin{equation}
    k_0 = \sigma_0\,\rho
    \label{eq:abs_coeff}
\end{equation}
For a line broadened only by the Doppler effect, the maximum absorption coefficient is :
\begin{equation}
    k_0 = \frac{\lambda_0^2}{8\pi}\cdot\frac{g_2}{g_1}\cdot\frac{\rho}{\Delta\nu_D\,\tau}
    \label{eq:k0}
\end{equation}
where $\lambda_0$ is the wavelength at the center of the absorption line, $\Delta\nu_D$ is
the Doppler width, $g_1$ and $g_2$ are the statistical weights of the lower and upper
states respectively, and $\tau$ is the radiative lifetime of the upper electronic state.
The Doppler width can be estimated from :
\begin{equation}
    \Delta\nu_D = \nu_0 \times 3\times10^{-7}
                  \left(\frac{T}{M}\right)^{1/2}
    \label{eq:doppler_width}
\end{equation}
where $\nu_0$ is the transition frequency, $T$ is the absolute temperature of the absorbing
gas, and $M$ is the atomic mass.
 
\subsubsection*{Selection Rules with Hyperfine Structure}
 
When hyperfine structure is taken into account, additional selection rules must be applied
for the total angular momentum quantum number $F$ and its projection $M$:
\begin{equation}
    \Delta F = 0,\,\pm 1,
    \qquad
    \Delta M = 0,\,\pm 1
    \label{eq:hfs_selection}
\end{equation}
In the optical pumping experiment, the incident light propagates parallel to the applied
magnetic field and is circularly polarized, so that only transitions with $\Delta M = +1$
(or $\Delta M = -1$) are allowed.
 
\subsubsection*{Magnetic Dipole Transitions}
 
We must also consider magnetic dipole transitions, which are approximately $10^5$ times
weaker than electric dipole transitions. These occur within the hyperfine structure and
between the magnetic sublevels, and are governed by :
\begin{equation}
    \Delta F = 0,\,\pm 1,
    \qquad
    \Delta M = 0,\,\pm 1
    \label{eq:mag_dipole_selection}
\end{equation}
In our experiment, the RF magnetic field is perpendicular to the dc magnetic field, so only
transitions with $\Delta F = 0$ and $\Delta M = \pm 1$ are observable. The $\Delta F = \pm 1$
transitions occur at RF frequencies of several gigahertz and cannot be observed with this
apparatus.
 
For allowed electric dipole transitions in emission, the lifetimes of the excited states are
of the order of $10^{-8}$~s, resulting in a natural linewidth of several hundred MHz. The
actual linewidth, determined by Doppler broadening, is of the order of 1~GHz. For
magnetic dipole transitions in the hyperfine structure of the ground electronic state the
radiative lifetimes are much longer, and collision processes determine the actual linewidths.
 